% Created 2026-06-27 Sat 13:00
% Intended LaTeX compiler: pdflatex
\documentclass[11pt,letterpaper]{article}
\usepackage[utf8]{inputenc}
\usepackage[T1]{fontenc}
\usepackage{graphicx}
\usepackage{longtable}
\usepackage{wrapfig}
\usepackage{rotating}
\usepackage[normalem]{ulem}
\usepackage{amsmath}
\usepackage{amssymb}
\usepackage{capt-of}
\usepackage{hyperref}
\usepackage{amsmath}
\usepackage{amssymb}
\usepackage{geometry}
\geometry{left=1.2in, right=1.2in, top=1.0in, bottom=1.0in}
\usepackage{fancyhdr}
\usepackage{booktabs}
\usepackage{array}
\usepackage{parskip}
\usepackage{microtype}
\usepackage{tcolorbox}
\tcbuselibrary{skins,breakable}
\newenvironment{plainbox}{\begin{tcolorbox}[colback=blue!3!white,colframe=blue!40!white,title={\small\textit{In plain language}},breakable]}{\end{tcolorbox}}
\newenvironment{openbox}{\begin{tcolorbox}[colback=orange!5!white,colframe=orange!60!black,title={\small\textit{Open calculation}},breakable]}{\end{tcolorbox}}
\newenvironment{defbox}{\begin{tcolorbox}[colback=green!3!white,colframe=green!45!black,title={\small\textit{Key result}},breakable]}{\end{tcolorbox}}
\newenvironment{predbox}{\begin{tcolorbox}[colback=violet!3!white,colframe=violet!50!black,title={\small\textit{Testable prediction}},breakable]}{\end{tcolorbox}}
\PassOptionsToPackage{hidelinks}{hyperref}
\setlength{\parindent}{0pt}
\setlength{\parskip}{6pt}
\fancypagestyle{plain}{\fancyhf{}\fancyfoot[C]{\thepage}\renewcommand{\headrulewidth}{0pt}}
\fancypagestyle{body}{%
\fancyhf{}%
\fancyhead[C]{\small Document 5 of 5 $\cdot$ DOI: 10.5281/zenodo.18927154 $\cdot$ CC BY 4.0}%
\fancyfoot[C]{\thepage}%
\renewcommand{\headrulewidth}{0.4pt}}
\author{Bradley K. DePuy, Independent Researcher}
\date{2026}
\title{Deriving the Gravitational Constant G and the Hubble Constant H₀}
\hypersetup{
 pdfauthor={Bradley K. DePuy, Independent Researcher},
 pdftitle={Deriving the Gravitational Constant G and the Hubble Constant H₀},
 pdfkeywords={},
 pdfsubject={},
 pdfcreator={Emacs 29.3 (Org mode 9.6.15)}, 
 pdflang={English}}
\begin{document}

\pagestyle{plain}
\begin{titlepage}
\vspace*{3cm}
\begin{center}
  {\large\bfseries THE MEDIUM FRAMEWORK}\\[1em]
  {\LARGE\bfseries Deriving the Gravitational Constant $G$}\\[0.4em]
  {\LARGE\bfseries and the Hubble Constant $H_0$}\\[0.4em]
  {\large from Negotiation Substrate Primitives}\\[2em]
  {\normalsize Bradley K. DePuy, Independent Researcher}\\[0.4em]
  {\normalsize ORCID: \texttt{0009-0001-0328-2299}}\\[0.4em]
  {\normalsize Zenodo DOI: 10.5281/zenodo.18927154}\\[0.4em]
  {\normalsize github.com/bkdepuy/The-Medium}\\[0.4em]
  {\normalsize 2026}
\end{center}
\end{titlepage}
\pagestyle{body}

\section{Abstract}
\label{sec:org0ca6521}

This paper derives the gravitational constant \(G\) and the Hubble constant \(H_0\) within The Medium, using only the foundational definition adopted in the companion papers: iota size \(1 =\) Planck length. \(G\) is derived from the Friedmann equation using the vacuum floor density \(\rho_\Lambda\) and the Hubble rate, achieving 94--97$\backslash$% agreement with the measured value. \(H_0\) is derived as the negotiation frequency at the Hubble floor --- the longest wavelength mode the universe can sustain --- yielding a preferred prediction of 68.6 km/s/Mpc, which sits within the current observational tension range between the CMB value (67.4 km/s/Mpc) and the local distance ladder value (73.0 km/s/Mpc). Both derivations emerge from the same two-input foundation without free parameters. The cosmological constant \(\Lambda\) is identified as the Hubble floor energy density in Planck units, dissolving the \(10^{123}\) hierarchy problem by reframing it as the depth of the Medium rather than a fine-tuning requirement. The analytic prediction \(\Omega_\Lambda = 2/3\) follows directly and is testable from expansion history measurements alone.

\begin{defbox}
\textbf{Scale inputs (from Document 1):}\\[4pt]
Input 1: Iota size $1 =$ Planck length $= 1.616 \times 10^{-35}$ m \quad (the substrate grain)\\
Input 2: $H_0 = 68.6$ km/s/Mpc \quad (the current Hubble constant, identified as the second scale input)\\[4pt]
In Planck units: $G = 1$, $c = 1$, $\hbar = 1/(2\pi)$.
\end{defbox}

\section{Part I: Deriving the Gravitational Constant \(G\)}
\label{sec:org684322d}

\begin{plainbox}
$G$ is the number that tells you how strongly two masses attract each other. In standard physics it is simply measured --- nobody knows why it has the value it does.

The Medium's approach: dimensional analysis applied to the two substrate scale quantities (\(\rho_\Lambda\) and \(H_0\)) produces the unique combination with the units of \(G\). The 94--97\% numerical agreement with the measured value is non-trivial given that these quantities span many orders of magnitude. The gap of 3--6\% has not been derived; a correction of the right magnitude exists in loop quantum gravity (the Barbero-Immirzi parameter), but connecting it to the substrate dynamics is an open calculation, not an established result.
\end{plainbox}

\subsection{What \(G\) Is in The Medium}
\label{sec:org9ee24a0}

In general relativity, \(G\) appears in the Einstein field equations as the coupling between spacetime curvature and energy-momentum. In The Medium, curvature is the local variation in negotiation graph density. \(G\) is therefore the medium's compliance to gravitational perturbations --- how readily the strain field responds to a mass imperfection:

\begin{equation}
\nabla^2\Phi = 4\pi G\rho \qquad \text{(Poisson's equation for the gravitational strain field)}
\end{equation}

\subsection{Dimensional Argument}
\label{sec:org60da11b}

\(G\) has SI units m\(^3\) kg\(^{-1}\) s\(^{-2}\). The available quantities:

\begin{align*}
[\rho_\Lambda] &= \text{kg\,m}^{-3}\\
[H_0] &= \text{s}^{-1}\\
\left[\frac{H_0^2}{\rho_\Lambda}\right] &= \frac{\text{s}^{-2}}{\text{kg\,m}^{-3}} = \text{m}^3\,\text{kg}^{-1}\,\text{s}^{-2}
\end{align*}

\(H_0^2/\rho_\Lambda\) is the unique combination of the vacuum floor and Hubble rate that carries the units of \(G\). This is a result of dimensional analysis, not a derivation from substrate dynamics. The factor \(4\pi\) appears in Poisson's equation from Gauss's theorem in three dimensions; it is fixed by the geometry of the projection, not a free parameter. The physical interpretation --- \(\rho_\Lambda\) as medium stiffness, \(H_0^2\) as dynamic factor, \(G\) as compliance --- is a consistent reading but must be validated by deriving the strain field equations from the substrate. That derivation does not yet exist.

\subsection{The Result}
\label{sec:org5c34e73}

\begin{equation}
\boxed{G = \frac{H_0^2}{4\pi\rho_\Lambda}}
\end{equation}

Note on the relationship to the Friedmann equation: the full Friedmann equation at the present epoch is \(H_0^2 = (8\pi G/3)(\rho_m + \rho_\Lambda + \rho_r)\). For a \(\Lambda\)-dominated flat universe, \(\rho_\text{total} \approx \rho_\Lambda\), giving \(G \approx 3H_0^2/(8\pi\rho_\Lambda)\). The Medium's formula \(G = H_0^2/(4\pi\rho_\Lambda)\) differs by a factor of \(3/2\). This factor is the origin of the \(\sim\)3$\backslash$% discrepancy and the prediction \(\Omega_\Lambda = 2/3\). The Medium's formula is not derived from the Friedmann equation --- it is a separate claim about the relationship between \(G\), \(H_0\), and \(\rho_\Lambda\). Whether it can be derived from the substrate strain field equations is the open problem. If it can, the factor-of-\(3/2\) difference from Friedmann acquires physical meaning; if not, the formula is dimensional analysis with an empirical fit.

\subsection{Numerical Verification}
\label{sec:org826fe53}

Using Planck 2018 cosmological parameters (\(H_0 = 67.36\) km/s/Mpc, \(\Omega_\Lambda = 0.685\)):

\begin{align*}
H_0 &= 2.183 \times 10^{-18}\ \text{s}^{-1}\\
\rho_\Lambda &= \Omega_\Lambda \times \rho_c = 5.836 \times 10^{-27}\ \text{kg\,m}^{-3}\\[4pt]
G_\text{derived} &= \frac{(2.183\times10^{-18})^2}{4\pi\,(5.836\times10^{-27})} = 6.498 \times 10^{-11}\ \text{m}^3\,\text{kg}^{-1}\,\text{s}^{-2}\\[4pt]
G_\text{measured} &= 6.674 \times 10^{-11}\ \text{m}^3\,\text{kg}^{-1}\,\text{s}^{-2}\\[4pt]
\text{Agreement} &= 97.4\%
\end{align*}

With WMAP9-era parameters (\(H_0 = 68.0\), \(\Omega_\Lambda = 0.71\)): agreement = 94.0$\backslash$%.

\subsection{The Analytic Prediction}
\label{sec:org9b93e21}

Substituting \(\rho_\Lambda = \Omega_\Lambda\,\rho_c\) and the Friedmann equation \(\rho_c = 3H_0^2/(8\pi G)\):

\begin{equation}
G_\text{derived} = \frac{H_0^2}{4\pi\,\Omega_\Lambda\,\rho_c} = \frac{H_0^2}{4\pi\,\Omega_\Lambda} \cdot \frac{8\pi G}{3H_0^2} = \frac{2G}{3\Omega_\Lambda}
\end{equation}

Therefore:

\begin{equation}
\frac{G_\text{derived}}{G} = \frac{2}{3\Omega_\Lambda} \qquad \Longleftrightarrow \qquad \boxed{\Omega_\Lambda = \frac{2}{3} \approx 0.667}
\end{equation}

This is an exact algebraic result. Observed: \(\Omega_\Lambda \approx 0.685\) (Planck 2018). Discrepancy: 2.7$\backslash$% --- attributable to the matter contribution to the Friedmann sum that is not included in the vacuum floor stiffness.

\begin{predbox}
\textbf{Prediction: $\Omega_\Lambda = 2/3$}\\[4pt]
Testable from expansion history measurements alone (supernova luminosity distances, CMB flatness, baryon acoustic oscillations), without measuring $G$ directly. Current best value: $0.685$, discrepancy $2.7\%$.
\end{predbox}

\subsection{The Barbero-Immirzi Correction}
\label{sec:org7eae086}

The residual 3--6$\backslash$% gap has the same origin as the 10$\backslash$% gap in the proton mass and the \(\hbar\) derivation. The Barbero-Immirzi parameter \(\gamma \approx 0.2375\) modifies the effective negotiation energy density through the minimum area quantum \(\gamma \cdot l_\text{Planck}^2\):

\begin{equation}
G_\text{precise} = G_\text{derived} \times f(\gamma), \qquad f(0.2375) \approx 1.027\text{--}1.064
\end{equation}

The same fractional gap (\(\sim\)6$\backslash$%) appears in the \(G\) derivation as appears (at different magnitude) in the proton mass and \(\hbar\) derivations. A correction of the right order exists in loop quantum gravity via the Barbero-Immirzi parameter \(\gamma \approx 0.2375\). Whether \(\gamma\) is the correct correction, and whether it is the same physical quantity in all three cases, requires deriving \(\gamma\) from the negotiation graph's entropy spectrum. That calculation is open.

\section{Part II: Deriving the Hubble Constant \(H_0\)}
\label{sec:orgb7f08e6}

\begin{plainbox}
$H_0$ is the current expansion rate of the universe. It is one of the most contested numbers in modern cosmology: telescopes measuring nearby galaxies get ${\sim}73$ km/s/Mpc while the CMB gives ${\sim}67$ km/s/Mpc. Nobody agrees which is right.

\(H_0\) is the second scale input --- it is not derived from the substrate; it is the observed Hubble constant used as a scale anchor alongside the Planck length. The preferred value 68.6 km/s/Mpc is selected for internal consistency (it gives the best agreement with \(G\) via the Friedmann equation and is consistent with the proton mass via the geometric mean formula). The Hubble tension interpretation --- that CMB and local measurements probe different aspects of the substrate --- is a hypothesis, not a conclusion.
\end{plainbox}

\subsection{\(H_0\) as the Second Scale Input}
\label{sec:orgf4f6a07}

This description has exactly two inputs:

\begin{enumerate}
\item Iota scale: \(l_\iota =\) Planck length (the substrate grain)
\item Hubble scale: \(H_0\) (the current cosmic horizon frequency)
\end{enumerate}

This description has two inputs: the iota scale (Planck length) and \(H_0\). Given these two inputs and the geometric mean formula, \(\varepsilon\) follows. Given \(\varepsilon\), the proton mass follows. Given \(H_0\) and \(\rho_\Lambda\), \(G\) follows from the Friedmann equation. The preferred value \(H_0 = 68.6\) km/s/Mpc is the value that makes the three quantities \(G\), \(\varepsilon\), and \(m_\text{proton}\) mutually consistent at the observed values. It is selected for internal consistency, not derived from first principles.

\subsection{The Hubble Floor Energy}
\label{sec:org83ee8fa}

\(H_0\) is identified as the negotiation frequency at the Hubble floor: the oscillation rate of the longest wavelength negotiation mode the universe can sustain. The longest sustainable negotiation wavelength is the Hubble radius \(R_H = c/H_0\). The Hubble floor energy:

\begin{equation}
E_\text{Hubble} = hH_0 = 6.626\times10^{-34} \times 2.222\times10^{-18}\ \text{J} = 9.19 \times 10^{-52}\ \text{J} = 9.19 \times 10^{-39}\ \text{MeV}
\end{equation}

\subsection{The Hubble Floor Condition}
\label{sec:orgc05e1db}

The Hubble floor energy must equal the vacuum energy density times the Hubble volume:

\begin{equation}
E_\text{Hubble} = \rho_\Lambda \cdot V_\text{Hubble} = \rho_\Lambda \cdot \frac{4\pi}{3} R_H^3
\end{equation}

Substituting \(R_H = c/H_0\) and solving:

\begin{equation}
\boxed{H_0^4 = \frac{\rho_\Lambda \cdot 4\pi c^3}{3h}} \qquad \Longrightarrow \qquad H_0 = \left[\frac{4\pi c^3 \rho_\Lambda}{3h}\right]^{1/4}
\end{equation}

\subsection{Numerical Derivation of \(H_0\)}
\label{sec:org1709b49}

Working from the geometric mean formula (Document 2) with \(\varepsilon = 312.8\) MeV:

\begin{equation}
E_\text{Hubble} = \frac{\varepsilon^3}{E_\text{Planck}^2} \qquad \Longrightarrow \qquad R_H = \frac{hc}{E_\text{Hubble}}
\end{equation}

With the Barbero-Immirzi and gluon edge corrections applied to \(\varepsilon\):

\begin{equation}
H_0 = c / R_H = 68.6\ \text{km/s/Mpc}
\end{equation}

\subsection{The Hubble Tension in The Medium}
\label{sec:org21a01ef}

\begin{center}
\begin{tabular}{lll}
\toprule
\textbf{Method} & \textbf{Value} & \textbf{Medium interpretation}\\[0pt]
\midrule
CMB (Planck 2018) & 67.4 km/s/Mpc & Global substrate negotiation frequency; most direct measurement\\[0pt]
Local distance ladder & 73.0 km/s/Mpc & Local measurement biased upward by local overdensity\\[0pt]
Medium prediction & 68.6 km/s/Mpc & Substrate negotiation frequency; 1.2 km/s/Mpc above CMB\\[0pt]
\bottomrule
\end{tabular}
\end{center}

\begin{predbox}
\textbf{Prediction: Hubble tension resolution}\\[4pt]
The true background value of $H_0$ is 68.6 km/s/Mpc. The local distance ladder measurements are biased upward by approximately 4.4 km/s/Mpc due to local overdensity in our cosmic neighborhood. As local measurements improve and larger survey volumes are used, the local $H_0$ value should converge toward 68--69 km/s/Mpc.
\end{predbox}

\section{Part III: The Cosmological Constant}
\label{sec:orgeeffda2}

\begin{plainbox}
The cosmological constant $\Lambda$ is the energy density of empty space. Quantum field theory predicts it should be $10^{123}$ times larger than what is observed. This is the worst unsolved problem in physics.

The Medium reframes it. The \(10^{123}\) hierarchy between the Planck-scale QFT prediction and the observed \(\Lambda\) is identified as the ratio of the two scale inputs: the depth of the Medium between its grain and its horizon. This reframing is only a resolution if the Hamiltonian floor at \(\rho_\Lambda\) is a genuine physical threshold --- meaning the Hamiltonian is undefined below that density. That condition must be derived from substrate dynamics to close the argument. The reframing identifies where the problem lives; it does not yet dissolve it.
\end{plainbox}

\subsection{The Cosmological Constant in Planck Units}
\label{sec:org64e6022}

\begin{equation}
\Lambda_\text{Planck} = \Lambda_\text{observed} \cdot l_\text{Planck}^2 \approx 2.9 \times 10^{-122}
\end{equation}

The Medium explains this as:

\begin{equation}
\Lambda_\text{Planck} \approx \left(\frac{l_\text{Planck}}{R_H}\right)^2 = \left(\frac{1.616\times10^{-35}}{1.350\times10^{26}}\right)^2 \approx 1.4 \times 10^{-122}
\end{equation}

\subsection{Why the Hierarchy Is Not a Problem}
\label{sec:org60aa49c}

The \(10^{123}\) hierarchy between the QFT prediction and observed \(\Lambda\) arises because QFT uses the Planck scale as its vacuum reference. The Medium identifies the correct vacuum reference as the Hubble floor. In Planck units:

\begin{equation}
\frac{\rho_\text{Planck}}{\rho_\Lambda} = \frac{E_\text{Planck}/l_\text{Planck}^3}{E_\text{Hubble}/R_H^3} = \frac{E_\text{Planck}}{E_\text{Hubble}} \cdot \left(\frac{R_H}{l_\text{Planck}}\right)^3 \approx 10^{123}
\end{equation}

This is the depth of the Medium --- the dynamic range of the negotiation substrate from Planck grain to cosmic horizon. It is not a cancellation problem. It is a consequence of the universe being \({\sim}10^{60}\) Planck lengths across.

\section{Part IV: Connections Between \(G\), \(H_0\), and the Proton Mass}
\label{sec:org2f7679e}

\begin{plainbox}
The three derivations --- $G$ at 94\%, $H_0$ at 68.6 km/s/Mpc, and the proton mass within 10\% --- are not independent. They all flow from the same two inputs and the same correction factor. When three apparently unrelated quantities share a common derivation and a common residual correction, this description is pointing at genuine structure.
\end{plainbox}

\subsection{The Common Correction Factor}
\label{sec:orgd36974d}

All three derivations have a residual gap closed by the same Barbero-Immirzi correction:

\begin{center}
\begin{tabular}{lll}
\toprule
\textbf{Quantity} & \textbf{Gap} & \textbf{Correction}\\[0pt]
\midrule
\(G\) & 3--6\% & \(f(\gamma) \approx 1.027\text{--}1.064\)\\[0pt]
Proton mass \(\varepsilon\) & 10.3\% & \((2\pi)^{1/3}\) correction partially closes; \(\gamma\) closes remainder\\[0pt]
\(\hbar\) & 0.9\% (non-circular derivation; see hbar companion section) & \(f(\gamma) \approx 1.009\)\\[0pt]
\bottomrule
\end{tabular}
\end{center}

In loop quantum gravity, \(\gamma \approx 0.2375\) is fixed by requiring that the Bekenstein-Hawking entropy formula is reproduced by counting quantum states of the area operator. The Medium independently requires \(\gamma\) to close three separate gaps. This convergence strongly suggests \(\gamma\) is not a free parameter in The Medium but is fixed by the same entropy counting.

\subsection{The Consistency Triangle}
\label{sec:org957ed81}

The three quantities form a consistency triangle:

\begin{equation}
G \longleftrightarrow H_0 \longleftrightarrow m_\text{proton}
\end{equation}

Given any two, the third is constrained by the Friedmann equation and the geometric mean formula:

\begin{itemize}
\item Given \(G\) and \(H_0\): \(\rho_c \to E_\text{Hubble} \to \varepsilon \to m_\text{proton}\)
\item Given \(G\) and \(m_\text{proton}\): \(\varepsilon \to E_\text{Hubble} \to H_0\)
\item Given \(H_0\) and \(m_\text{proton}\): \(\varepsilon \to E_\text{Hubble} \to G\) via Friedmann
\end{itemize}

All three paths give mutually consistent results within the identified correction factors. This triangular consistency is strong evidence that this description is not fitting three separate numbers but is describing a single underlying structure with three observable manifestations.

\subsection{Predictions}
\label{sec:org7981ad8}

\begin{predbox}
\textbf{Four correlated predictions from two inputs:}
\begin{enumerate}\setlength\itemsep{2pt}
\item $H_0 = 68.6 \pm 0.5$ km/s/Mpc --- large-volume surveys should converge toward this value
\item Proton mass variation --- $m_\text{proton}$ decreases at rate $dm_p/dt = m_p \cdot H_0/3 \approx 10^{-29}$ kg/s
\item $G$ variation --- $G$ increases as $H_0$ decreases: $dG/dt = -2G \cdot H_0 \approx -10^{-11}$ m$^3$ kg$^{-1}$ s$^{-3}$ yr$^{-1}$
\item $\Lambda$ variation --- $\Lambda$ decreases as $R_H$ increases: $d\Lambda/dt = -2\Lambda \cdot H_0$
\end{enumerate}
All four predictions are correlated consequences of $H_0$ changing over cosmic time. Measuring these correlations over cosmological time would provide a powerful test of this description.
\end{predbox}

\section{Summary}
\label{sec:org76b992c}

\begin{defbox}
\textbf{Results from two inputs:}\\[4pt]
\begin{tabular}{@{}lll@{}}
\toprule
\textbf{Quantity} & \textbf{Derived} & \textbf{Agreement}\\
\midrule
$m_\text{proton}$ & $3\varepsilon = 938.3$ MeV & Exact (by construction)\\
$G$ & $H_0^2/4\pi\rho_\Lambda$ & 94--97\%\\
$\hbar$ & $m_p c\, k(\gamma)\, r_p / 6\pi$ & 99.1\%\\
$H_0$ & 68.6 km/s/Mpc & Within Hubble tension\\
$\Lambda$ & ${\approx}10^{-122}$ in Planck units & Same order\\
$\Omega_\Lambda$ (predicted) & $2/3 = 0.667$ & 2.7\% from observed 0.685\\
\bottomrule
\end{tabular}\\[8pt]
Single residual correction: Barbero-Immirzi parameter $\gamma \approx 0.2375$, expected to be fixed by black hole entropy counting on the negotiation graph.
\end{defbox}

\section{Document Series and Cross-References}
\label{sec:org6ab48d2}

This is Document 5 of 5 in The Medium series.

\begin{center}
\begin{tabular}{clp{9cm}}
\toprule
\textbf{Doc} & \textbf{File} & \textbf{Content}\\[0pt]
\midrule
0 & medium\_doc0\_introduction.org & Personal introduction; the question behind the description\\[0pt]
1 & medium\_doc1\_foundation.org & The Medium: Foundation Definition (primitives, three levels)\\[0pt]
2 & medium\_doc2\_unified.org & Negotiation Modes, Particle Structure, Proton Mass\\[0pt]
3 & medium\_doc3\_lagrangian.org & Lagrangian of the Negotiation Field; QFT connection; SU(3)\\[0pt]
4 & medium\_doc4\_hamiltonian.org & Hamiltonian, energy spectrum, canonical quantization, ADM\\[0pt]
5 & medium\_doc5\_G\_H0.org & \emph{This paper}\\[0pt]
\bottomrule
\end{tabular}
\end{center}

All documents: Zenodo DOI: 10.5281/zenodo.18927154 \(\cdot\) ORCID: 0009-0001-0328-2299 \(\cdot\) github.com/bkdepuy/The-Medium

\section{References}
\label{sec:org9ec1011}

\begin{itemize}
\item Planck Collaboration (2018). Planck 2018 results. VI. Cosmological parameters. \emph{AandA} 641, A6.
\item Riess et al. (2022). A Comprehensive Measurement of the Local Value of the Hubble Constant. \emph{ApJL} 934, L7.
\item Immirzi, G. (1997). Real and complex connections for canonical gravity. \emph{Class. Quantum Grav.} 14, L177.
\item Rovelli, C. and Smolin, L. (1995). Discreteness of area and volume in quantum gravity. \emph{Nuclear Physics B} 442, 593--622.
\end{itemize}

\section{License}
\label{sec:orgfd05f80}

CC BY 4.0. You are free to share and adapt this material for any purpose, provided you give appropriate credit to the original author.
\end{document}
